\documentclass[11pt,a4paper]{article}

\usepackage{amsmath,amssymb,amsthm}
\usepackage[margin=1in]{geometry}
\usepackage{bm}
\usepackage{hyperref}
\usepackage{xcolor}

\newcommand{\Jc}{\mathcal{J}}
\newcommand{\Hc}{\mathcal{H}}
\newcommand{\Jh}{\hat J}
\newcommand{\kB}{k_{\mathrm B}}
\newcommand{\bk}{\vec k}
\newcommand{\bQ}{\vec Q}
\newcommand{\bq}{\vec q}
\newcommand{\brr}{\vec r}
\newcommand{\bdel}{\vec\delta}
\newcommand{\dk}{\frac{d^3k}{(2\pi)^3}}
\newcommand{\avg}[1]{\left\langle #1 \right\rangle}
\newcommand{\Wc}{\mathcal{W}}
\newcommand{\TODO}[1]{\textcolor{red}{\textbf{[TODO: #1]}}}
\newcommand{\FLAG}[1]{\textcolor{blue}{\textbf{[#1]}}}

\title{Notes: the planar antiferromagnetic mean spherical model\\
as a thermodynamic-geometry example}
\author{}
\date{Working notes --- \S0--\S2, second draft}

\begin{document}
\maketitle

\noindent
These notes develop the $d=3$ planar antiferromagnetic \emph{mean} spherical model (MSM)
as a replacement for the mean-field antiferromagnetic Ising example in the manuscript.
Reference [K] is Knops, Phys.\ Rev.\ B \textbf{8}, 4209 (1973); equation numbers of the
form (K2.39) refer to that paper. Reference [BR] is the current draft (Basri--Raz).

Everything is written out. Where a factor of $2$ or a $\tfrac12$ appears it is derived,
not asserted, because the two papers use different conventions and the first draft of
these notes got several of them wrong.

Two decisions are fixed: the mean spherical ensemble, and $H'=0$ in the Hamiltonian with
control manifold $(\beta,\beta\Hc)$.

%===============================================================
\section{Conventions}
%===============================================================

\subsection{Lattice}
Simple cubic, $N=L^3$ sites, unit lattice constant, periodic boundary conditions.
It is bipartite: $\brr_j\in A$ if $\sum_\mu (r_j)_\mu$ is even, $\in B$ if odd.
Each sublattice has $N/2$ sites. Nearest neighbours of an $A$ site all lie in $B$.

\subsection{Rescaled couplings: why we do not use $J$ and $H$ directly}
\label{sec:rescale}
Knops writes the planar energy (K4.1) with explicit factors of $2$ on both the coupling
and the field. Those factors are not decoration: the planar constraint (K4.2) is
$\avg{|\bm\sigma_j|^2}=1$ summed over \emph{two} components, so $\avg{(\sigma_j^z)^2}=\tfrac12$
rather than $1$, and doubling $J$ and $H$ is exactly what compensates. That is the entire
content of his claim that the planar model has the same phase boundary and the same
longitudinal averages as the scalar model.

Carrying those $2$'s around is a reliable source of error --- in particular the control
parameter conjugate to the magnetisation is $2\beta H$, not $\beta H$, and getting this
wrong misstates every field-direction component of the metric by a factor of $4$. We
therefore absorb them once and for all:
\begin{equation}
  \boxed{\ \Jc \equiv 2J, \qquad \Hc \equiv 2H\ }
\end{equation}
and work exclusively with $\Jc,\Hc$. A dictionary back to [K] is given in \S\ref{sec:dict}.

The energy is then, with a \emph{bond} sum (each unordered nearest-neighbour pair once),
\begin{equation}
  E \;=\; \sum_{\langle ij\rangle} \Jc_{ij}\,\bm\sigma_i\cdot\bm\sigma_j
      \;-\; \Hc\sum_j \sigma_j^z ,
  \qquad
  \Jc_{ij}=\Jc>0 \ \text{ for $i\in A$, $j\in B$ n.n.}
  \label{eq:energy}
\end{equation}
Following [K] the sign convention is the opposite of the usual one: \emph{positive} $\Jc$
is antiferromagnetic. This must be flagged in the manuscript, which uses $K=-J$ with $J<0$.

\subsection{Fourier transform}
Symmetric convention, \textbf{full} first Brillouin zone:
\begin{equation}
  \bm\sigma_j = \frac{1}{\sqrt N}\sum_{\bk} e^{i\bk\cdot\brr_j}\,\bm S_{\bk},
  \qquad
  \bm S_{\bk} = \frac{1}{\sqrt N}\sum_j e^{-i\bk\cdot\brr_j}\,\bm\sigma_j ,
  \label{eq:ft}
\end{equation}
with $\bk=(2\pi/L)(l_x,l_y,l_z)$, $l_\mu=0,\dots,L-1$, so there are exactly $N$ wavevectors.
The basic identities are
\begin{equation}
  \sum_j e^{i(\bk-\bk')\cdot\brr_j} = N\,\delta_{\bk\bk'},
  \qquad
  \bm S_{-\bk} = \bm S_{\bk}^{*},
  \qquad
  \sum_j |\bm\sigma_j|^2 = \sum_{\bk} |\bm S_{\bk}|^2 ,
  \label{eq:parseval}
\end{equation}
the second from reality of $\bm\sigma_j$, the third being Parseval. Throughout
$|\bm S_{\bk}|^2 \equiv |S^x_{\bk}|^2+|S^z_{\bk}|^2$.

We use the full zone, not the half zone of [K]. Knops restricts to half the zone
(remark below K2.12) because his two-sublattice formulation double counts; with $H'=0$
that formulation collapses (\S\ref{sec:collapse}) and the restriction is unnecessary.
It is also where his factors of $2$ become hardest to follow.

\subsection{The structure function}
Define
\begin{equation}
  \Jc(\bk) \;\equiv\; \tfrac12\sum_{\bdel} e^{-i\bk\cdot\bdel}\,\Jc_{\bdel}
  \;=\; \Jc\left(\cos k_x+\cos k_y+\cos k_z\right),
  \qquad \Jc(0)=3\Jc ,
  \label{eq:structure}
\end{equation}
the second equality because the six nearest-neighbour vectors $\bdel=\pm\hat e_\mu$
contribute $\tfrac12\Jc\sum_\mu(e^{-ik_\mu}+e^{ik_\mu})$.

\paragraph{Why the $\tfrac12$.} It is fixed by requiring that the bond sum diagonalise with
unit coefficient. Writing the bond sum as half the unrestricted sum over ordered pairs,
\begin{equation}
  \sum_{\langle ij\rangle}\Jc_{ij}\,\bm\sigma_i\cdot\bm\sigma_j
  = \tfrac12\sum_j\sum_{\bdel}\Jc_{\bdel}\,\bm\sigma_j\cdot\bm\sigma_{j+\bdel},
\end{equation}
and inserting \eqref{eq:ft},
\begin{align}
  \sum_j \bm\sigma_j\cdot\bm\sigma_{j+\bdel}
  &= \frac{1}{N}\sum_j\sum_{\bk,\bk'} e^{i\bk\cdot\brr_j}\,e^{i\bk'\cdot(\brr_j+\bdel)}\,
     \bm S_{\bk}\cdot\bm S_{\bk'} \nonumber\\
  &= \frac{1}{N}\sum_{\bk,\bk'} e^{i\bk'\cdot\bdel}\,\bm S_{\bk}\cdot\bm S_{\bk'}\;
     N\delta_{\bk',-\bk}
   \;=\; \sum_{\bk} e^{-i\bk\cdot\bdel}\,|\bm S_{\bk}|^2 ,
\end{align}
using \eqref{eq:parseval} twice. Hence
\begin{equation}
  \boxed{\ \sum_{\langle ij\rangle}\Jc_{ij}\,\bm\sigma_i\cdot\bm\sigma_j
   = \sum_{\bk}\Jc(\bk)\,|\bm S_{\bk}|^2\ }
  \label{eq:diag}
\end{equation}
with no residual factor. This is the same $\tfrac12$ that appears in (K2.10), and reading
Knops' $\sum_{\vec i,\vec j}$ as an unrestricted double sum instead of a bond sum would
give $2\Jc(\bk)$ here and $\kB T_0 = 4/W_3$ instead of $2/W_3$ --- which is how one confirms
the bond reading is the intended one.

The field term is likewise
\begin{equation}
  \sum_j \sigma_j^z = \sqrt N\, S_0^z .
  \label{eq:fieldterm}
\end{equation}

\subsection{Bipartite identity and the ordering wavevector}
Since $\cos(\pi-k_\mu)=-\cos k_\mu$,
\begin{equation}
  \Jc(\bQ-\bk) = -\Jc(\bk), \qquad \bQ\equiv(\pi,\pi,\pi),
  \label{eq:bipartite}
\end{equation}
which is (K2.11). In particular $\Jc(\bQ)=-\Jc(0)$, and since $\Jc>0$ the minimum of
$\Jc(\bk)$ is at $\bk=\bQ$: the ordering wavevector is $\bQ$, i.e.\ two-sublattice
antiferromagnetic order. Note $2\bQ$ is a reciprocal lattice vector, so $\bQ\equiv-\bQ$
and $\bm S_{\bQ}$ is \emph{real}.

Near the zone corner, writing $\bk=\bQ-\bq$,
\begin{equation}
  \Jc(\bQ-\bq) = -\Jc(\bq) = -\Jc\sum_\mu\cos q_\mu
   = -\Jc(0) + \tfrac{\Jc}{2}q^2 + O(q^4).
  \label{eq:smallq}
\end{equation}

\subsection{The symbol $\beta$, and the control parameters}
Three-way collision ($\beta$ = inverse temperature, order-parameter exponent, and inside
$b=2-\alpha$). Here $\beta\equiv1/\kB T$ always; the order-parameter exponent is written
$\beta_{\rm op}$.

The Boltzmann weight is $\exp(-\beta E)$ with \eqref{eq:energy}, i.e.
\begin{equation}
  \exp\Big(-\beta E_{\rm int} + h\sum_j\sigma_j^z\Big),
  \qquad
  E_{\rm int}\equiv\sum_{\langle ij\rangle}\Jc_{ij}\bm\sigma_i\cdot\bm\sigma_j,
  \qquad
  h \equiv \beta\Hc .
\end{equation}
This is manifestly linear in
\begin{equation}
  \boxed{\ \lambda=(\lambda^1,\lambda^2)=(\beta,\,h),\qquad h=\beta\Hc=2\beta H\ }
  \label{eq:controls}
\end{equation}
as required by [BR, Eq.~(1)], with conjugate variables
\begin{equation}
  X_1=-E_{\rm int},\qquad X_2=M_{\rm tot}=\sum_j\sigma^z_j ,
  \qquad
  \frac{\partial\psi}{\partial\lambda^\mu}=\avg{X_\mu},
  \qquad \psi\equiv\ln Z .
\end{equation}

\subsection{Dictionary to [K]}
\label{sec:dict}
\begin{center}
\begin{tabular}{lll}
\hline
these notes & [K] & relation \\
\hline
$\Jc$ & $J$ & $\Jc=2J$ \\
$\Jc(\bk)$ & $\Jh(\bk)$ & $\Jc(\bk)=2\Jh(\bk)$, \ $\Jc(0)=2\Jh(0)=6J$ \\
$\Hc$ & $H$ & $\Hc=2H$ \\
$M=\avg{\sigma^z}$ & $M$ & identical \\
$x^2$ & $x^2$ & identical (K2.35, K4.8) \\
$\Wc(z)$ & $W_3$ & $\Wc(z)=\tfrac{1}{2\beta}W_3\!\big(z/2\beta\big)$ \\
\hline
\end{tabular}
\end{center}

\subsection{Numerical constants (nearest-neighbour sc)}
The Watson integral is
\begin{equation}
  W_{\rm sc}\equiv\int_{\rm BZ}\!\dk\;
   \frac{1}{1-\tfrac13(\cos k_x+\cos k_y+\cos k_z)} = 1.5163860592\ldots
\end{equation}
As derived in \S\ref{sec:T0},
\begin{equation}
  \kB T_0=\frac{3\Jc}{W_{\rm sc}}=\frac{6J}{W_{\rm sc}}=3.9568\,J,
  \qquad
  \beta_0\Jc(0)=W_{\rm sc}=1.5164,
  \qquad
  \Hc_c(T{=}0)=4\Jc(0)=24J .
\end{equation}
The first reproduces $\kB T_0/J=3.957$ quoted below (K2.29); the last is
$H_c(0)=12J=4\Jh(0)$ in Knops' variables, twice the Ising value as he notes in \S III.

%===============================================================
\section{The model}
%===============================================================

\subsection{Spins and constraints}
Each site carries a two-component (planar) real spin $\bm\sigma_j=(\sigma_j^x,\sigma_j^z)$,
subject to spherical constraints imposed separately on each sublattice (K4.2):
\begin{equation}
  \sum_{j\in A}|\bm\sigma_j|^2=\tfrac12 N,
  \qquad
  \sum_{j\in B}|\bm\sigma_j|^2=\tfrac12 N .
  \label{eq:constraints}
\end{equation}
Two constraints rather than one, because a staggered field breaks translational invariance
in general (discussion below K2.5). In the mean spherical version these hold on average,
enforced by two spherical fields $z_A,z_B$ (K2.6).

\subsection{Why planar rather than scalar}
The scalar antiferromagnetic SM does have long-range order --- the correlation functions
approach a constant $\propto x^2$ (K2.35) --- but at $H\neq0$ that order is invisible in
every one-point function: $\lim_{H'\to0}M'=0$ (below K2.43). The order parameter has no
thermodynamic conjugate. For thermodynamic geometry that is fatal, since no control
parameter couples to it at all.

The planar model repairs this. With spins in a plane and a transverse staggered field
$H'_\perp$, [K,\S IV] shows (K4.7--4.8)
\begin{equation}
  \lim_{H'_\perp\to0}(S^\perp_A)^2=\frac{x^2}{\lambda},
  \qquad
  \lim_{H'_\perp\to0}(S^\perp_B)^2=x^2\lambda ,
\end{equation}
so $x^2$ \emph{is} the transverse spontaneous staggered magnetisation. The model then
mimics the isotropic $n$-vector antiferromagnet, spin-flop physics included, while keeping
exact solubility and --- by the rescaling of \S\ref{sec:rescale} --- the identical phase
boundary and longitudinal averages of the scalar case. We inherit all of [K,\S II] and
gain an order parameter.

\subsection{Mean versus strict spherical: what actually differs}
\label{sec:ensembles}
The two ensembles are equivalent for thermodynamics. The Massieu potential per site is the
same, hence so is every derivative of it: susceptibilities, specific heat, and the full
static Hessian $\partial^2\psi/\partial\lambda^\mu\partial\lambda^\nu$ are
ensemble-independent. Essentially all of \S2 is therefore insensitive to the choice.

The real distinction is not statics versus dynamics but \emph{thermodynamic derivatives
versus fluctuation formulas evaluated at fixed $z$}. In the MSM, $\sum_j|\bm\sigma_j|^2$
genuinely fluctuates, with
\begin{equation}
  \mathrm{Var}\Big(\sum_j|\bm\sigma_j|^2\Big)\ \sim\ N\!\int\! d^dq\; G(\bq)^2 ,
  \qquad G(\bq)\sim\frac{1}{q^2}\ \text{in the ordered phase},
\end{equation}
i.e.\ $\int d^dq/q^4$, convergent only for $d>4$. In $d=3$ the fluctuations of the
constrained quantity are therefore \emph{superextensive}. (Note this is not a statement
that the $z$-saddle is broad: $\Phi_{zz}\sim I_2\to\infty$ makes it \emph{narrower},
$\delta z\sim(NI_2)^{-1/2}$. The point is that conditioning on $Q$ changes fluctuation
formulas for extensive \emph{quadratic} observables.) Consequently
\begin{equation}
  \avg{\delta X_\mu\,\delta X_\nu}_{z}\ \neq\
  \frac{\partial^2\psi}{\partial\lambda^\mu\partial\lambda^\nu} .
  \label{eq:notequal}
\end{equation}
(No factor $\beta$: in the natural Massieu coordinates \eqref{eq:controls},
$\partial_\mu\partial_\nu\psi=\mathrm{Cov}(X_\mu,X_\nu)$ identically.)
This single integral is simultaneously Knops' observation that $\chi'_{\rm MSM}=\infty$
throughout the ordered region while $\chi'$ is finite (K2.36--2.38), and the cancellation
made explicit in \S\ref{sec:cancel} below. Knops was aware of it: his footnote 9 hangs off
the phrase ``when sufficient care is taken.''

\paragraph{Working rule.} Obtain static quantities by differentiating $\psi$; never by a
fixed-$z$ fluctuation formula. \S\ref{sec:ordered} shows this is nearly free here.

\FLAG{Consequence for \S3.} There is no such escape route for $T_{\mu\nu}$, which
\emph{is} a fluctuation object. A Langevin dynamics that lets $\sum_j|\bm\sigma_j|^2$ float
will equilibrate to the fixed-$z$ spectrum, i.e.\ the wrong one, and will inherit the
divergence of \eqref{eq:notequal} directly. The spherical field must be dynamical and
slaved. This is the one place where the ensemble choice imposes a real obligation.

\subsection{Why $H'=0$, and what it costs}
Setting $H'=0$ in the Hamiltonian buys: $z_A=z_B$ by symmetry (below K2.6), hence the
collapse of \S\ref{sec:collapse}; a two-dimensional control manifold matching [BR, Fig.~1];
and a closed-form critical line.

The cost is specific. Removing $H'_\perp$ removes the field conjugate to the order
parameter. Moreover $\bm\sigma\to(-\sigma^x,\sigma^z)$ is a symmetry of \eqref{eq:energy},
so
\begin{equation}
  \avg{\delta M_{\rm tot}\,\delta S_\perp}=\avg{\delta E_{\rm int}\,\delta S_\perp}=0 ,
  \label{eq:decoupling}
\end{equation}
and the order-parameter susceptibility \emph{never enters $g_{\mu\nu}$ at linear order}.
The metric is built entirely from energy-like and uniform-magnetisation-like fluctuations.
This is the reason the static Hessian turns out to be finite at the transition
(\S\ref{sec:crit}); it is a physical statement about the example --- one drives the system
through the transition with fields that couple only to non-critical one-point functions ---
not an artefact.

\FLAG{Corrected.} An earlier version of these notes claimed that the degeneracy of the
two components in the fixed-$z$ kernel implies a spontaneously broken $O(2)$ and a
Goldstone mode. That is wrong. The term $-\Hc\sum_j\sigma^z_j$ breaks $O(2)$ down to the
reflection $\sigma^x\to-\sigma^x$; the transverse staggered order breaks that residual
$Z_2$, so at $\Hc\neq0$ there is \emph{no} Goldstone mode. Only at $\Hc=0$ is $O(2)$ restored.

The mechanism that lifts the degeneracy is the \emph{second} spherical constraint, and it
is worth spelling out because it is the spin-flop selection. Write
$\sigma^z_j=M+\epsilon_j m_s$ with $\epsilon_j=\pm1$ on the two sublattices. Then
$|\bm\sigma_j|^2$ contains a term $2M\epsilon_j m_s$, so
\begin{equation}
  Q_-\equiv\sum_{j\in A}|\bm\sigma_j|^2-\sum_{j\in B}|\bm\sigma_j|^2 \;\ni\; 2NMm_s ,
\end{equation}
and a longitudinal staggered component violates $Q_-=0$ at linear order whenever $M\neq0$.
A transverse component contributes $\epsilon_j^2 x^2=x^2$ and does not. The constraint
therefore forbids longitudinal staggered order in a uniform field and selects the
transverse channel --- exactly the spin-flop physics of [K,\S IV]. Longitudinal and
transverse staggered fluctuations are \emph{not} degenerate.

\subsection{Collapse to a one-sublattice problem}
\label{sec:collapse}
Knops' quadratic form (K2.9) contains a sublattice-mixing term
$-\tfrac12(z_A-z_B)S_{\bk}S^*_{\bQ-\bk}$. With $H'=0$ we have $z_A=z_B\equiv z$ and
\emph{this term vanishes identically}. What remains is a translationally invariant mean
spherical model with a single spherical field, a single scalar propagator, and ordering at
$\bQ$. Combining \eqref{eq:diag} and \eqref{eq:fieldterm}, the exponent of the Boltzmann
weight including the spherical term $-z\sum_j|\bm\sigma_j|^2$ is
\begin{equation}
  \boxed{\ -\sum_{\bk} a_{\bk}\,|\bm S_{\bk}|^2 \;+\; h\sqrt N\,S^z_0 ,
  \qquad a_{\bk}\equiv z+\beta\Jc(\bk)\ }
  \label{eq:quadform}
\end{equation}
By \eqref{eq:bipartite} the soft mode is at $\bk=\bQ$, where $a_{\bQ}=z-\beta\Jc(0)$.

\FLAG{Scope of this collapse.} It holds for the \emph{equilibrium mean-spherical saddle},
and hence for all of \S2: there $z_A=z_B$ and $Q_-=0$ are satisfied automatically, and no
$2\times2$ sublattice structure is needed. It does \emph{not} hold for strictly constrained
dynamics. Enforcing $Q_+=N$ and $Q_-=0$ configuration by configuration requires two
time-dependent multipliers, and the $Q_-$ constraint couples $\bk$ to $\bQ-\bk$. That
coupling is precisely what distinguishes longitudinal from transverse staggered
fluctuations in a uniform field (\S1.4), so it cannot be dropped in \S3--\S4.

%===============================================================
\section{The Massieu potential}
%===============================================================

\subsection{Gaussian integration at fixed $z$}
It is cleanest to count degrees of freedom in real space. Write \eqref{eq:quadform} as
\begin{equation}
  -\sum_{ij}A_{ij}\big(\sigma^x_i\sigma^x_j+\sigma^z_i\sigma^z_j\big)+h\sum_j\sigma^z_j ,
  \qquad
  A_{ij}=z\,\delta_{ij}+\tfrac{\beta}{2}\Jc_{ij},
\end{equation}
so that $A$ is a real symmetric $N\times N$ matrix with eigenvalues $a_{\bk}$ (the
$\tfrac12$ because the unrestricted double sum is twice the bond sum). Both components see
the same $A$; only $\sigma^z$ carries a source, with $b_j=h$ for every $j$.

Using $\int d^Nu\,e^{-u^\top A u+b^\top u}=\pi^{N/2}(\det A)^{-1/2}e^{\,b^\top A^{-1}b/4}$
once per component,
\begin{equation}
  Z_z=\underbrace{\frac{\pi^{N/2}}{(\det A)^{1/2}}}_{x\text{-component}}
      \cdot
      \underbrace{\frac{\pi^{N/2}}{(\det A)^{1/2}}\,e^{\,b^\top A^{-1}b/4}}_{z\text{-component}}
     =\frac{\pi^{N}}{\det A}\,e^{\,b^\top A^{-1}b/4}.
\end{equation}
For the source term, the uniform vector $u_j=1$ has Fourier components $\sqrt N\delta_{\bk,0}$, so
\begin{equation}
  b^\top A^{-1}b = h^2\sum_{ij}(A^{-1})_{ij}=h^2\sum_{\bk}\frac{|\hat u_{\bk}|^2}{a_{\bk}}
   =\frac{h^2 N}{a_0}.
\end{equation}
With $\ln\det A=\sum_{\bk}\ln a_{\bk}$,
\begin{equation}
  \boxed{\ \ln Z_z=-\sum_{\bk}\ln a_{\bk}+\frac{h^2N}{4a_0}+N\ln\pi\ }
  \label{eq:lnZ}
\end{equation}

\FLAG{Correction to the first draft.} The first draft had $a_{\bk}=z+\beta\Jh(\bk)$ and
omitted the $\tfrac14$ in the source term. Both are fixed above.

\paragraph{Immediate check.} The source couples as $h\sum_j\sigma^z_j$, so
\begin{equation}
  N M \equiv \avg{\textstyle\sum_j\sigma^z_j}=\frac{\partial \ln Z_z}{\partial h}
  =\frac{hN}{2a_0}
  \qquad\Longrightarrow\qquad
  M=\frac{h}{2a_0},
  \label{eq:M}
\end{equation}
which is also what one gets directly from the single displaced mode
$-a_0|S^z_0|^2+h\sqrt N S^z_0$, whose mean is $\avg{S^z_0}=h\sqrt N/2a_0$.

\subsection{Restoring the constraint}
$Z_z$ is the partition function with the spherical term inserted by hand. The constrained
one follows by a Legendre/saddle-point step: the physical Massieu potential is
\begin{equation}
  \frac{\psi}{N}=\operatorname*{extr}_{z}\left[\,z+\frac{1}{N}\ln Z_z\right],
\end{equation}
since stationarity gives
$1-\tfrac1N\avg{\sum_j|\bm\sigma_j|^2}=0$, which is precisely \eqref{eq:constraints} in the
mean. Adding a macroscopic occupation $x$ of the soft mode --- put $|S^x_{\bQ}|^2=Nx^2$,
contributing $-a_{\bQ}Nx^2$ to \eqref{eq:quadform} --- gives the final form
\begin{equation}
  \boxed{\;
  \frac{\psi}{N}=\operatorname*{extr}_{z,\,x}\ \Phi,
  \qquad
  \Phi = z-\int\!\dk\,\ln\!\big[z+\beta\Jc(\bk)\big]
        +\frac{h^2}{4\big[z+\beta\Jc(0)\big]}
        -\big[z-\beta\Jc(0)\big]x^2 \;}
  \label{eq:Phi}
\end{equation}
up to an additive constant independent of $(\beta,h)$. Convergence of the Gaussian
integral requires $a_{\bk}>0$ for all $\bk$, i.e.
\begin{equation}
  z\ \geq\ \beta\Jc(0) ,
  \label{eq:zrange}
\end{equation}
the analogue of (K2.23).

\subsection{Stationarity: sticking and the constraint}
Differentiating \eqref{eq:Phi} in $x$:
\begin{equation}
  -2x\big[z-\beta\Jc(0)\big]=0
  \quad\Longrightarrow\quad
  \underbrace{x=0}_{\text{disordered}}
  \quad\text{or}\quad
  \underbrace{z=\beta\Jc(0)}_{\text{ordered}} .
  \label{eq:sticking}
\end{equation}
This is the standard spherical-model mechanism: $z$ \emph{sticks} at the lower end of its
range \eqref{eq:zrange} throughout the ordered region (below K2.23).

Differentiating in $z$, and using \eqref{eq:M} in the form $M^2=h^2/4a_0^2$:
\begin{equation}
  \boxed{\ \Wc(z)+M^2+x^2=1\ },
  \qquad
  \Wc(z)\equiv\int\!\dk\,\frac{1}{z+\beta\Jc(\bk)} .
  \label{eq:constraint}
\end{equation}
This is the $H'=0$ ($\lambda=1$) specialisation of (K2.27)--(K2.28), which then become
identical to one another. Using \eqref{eq:bipartite} and invariance of the zone average
under $\bk\to\bQ-\bk$,
\begin{equation}
  \Wc(z)=\int\!\dk\,\frac{1}{z-\beta\Jc(\bk)}=\frac{1}{2\beta}\,W_3\!\left(\frac{z}{2\beta}\right),
\end{equation}
with $W_3$ the integral (K2.19).

\paragraph{Disordered region.} $x=0$, and \eqref{eq:constraint} fixes $z>\beta\Jc(0)$
implicitly. Define
\begin{equation}
  r\equiv z-\beta\Jc(0)>0 ,
\end{equation}
so by \eqref{eq:smallq} $a_{\bQ-\bq}=r+\tfrac{\beta\Jc}{2}q^2$, giving $\xi^2=\beta\Jc/2r$
and $\nu=1$.

\paragraph{Ordered region.} $r=0$ identically, and \eqref{eq:constraint} fixes the
condensate. \FLAG{Note} that $r=0$ holds throughout the ordered phase, not merely on the
boundary: the whole ordered region is gapless. This is the origin of the divergence of
$\chi'_{\rm MSM}$ in \S\ref{sec:ensembles} and the main technical risk for \S3.

\subsection{The zero-field critical temperature}
\label{sec:T0}
Set $H=0$, $x=0$, $z=\beta_0\Jc(0)$ in \eqref{eq:constraint}, so $\Wc=1$:
\begin{equation}
  \frac{1}{\beta_0}\int\!\dk\,\frac{1}{\Jc(0)-\Jc(\bk)}=1 .
\end{equation}
With $\Jc(0)-\Jc(\bk)=3\Jc\big[1-\tfrac13\sum_\mu\cos k_\mu\big]$ this is
$W_{\rm sc}/(3\Jc\beta_0)=1$, hence
\begin{equation}
  \boxed{\ \kB T_0=\frac{3\Jc}{W_{\rm sc}}=\frac{6J}{W_{\rm sc}}=3.9568\,J\ },
  \qquad
  \beta_0\Jc(0)=W_{\rm sc}.
\end{equation}
Converting with $\Jc=2J$: $W_3(\Jh(0))=2/\kB T_0$, i.e.\ (K2.29), and $\kB T_0/J=3.957$. \checkmark

\subsection{The critical line}
On the boundary both $x=0$ and $z=\beta\Jc(0)$ hold, so $a_0=2\beta\Jc(0)$ and by
\eqref{eq:M}
\begin{equation}
  M_c=\frac{h}{4\beta\Jc(0)}=\frac{\Hc_c}{4\Jc(0)} .
\end{equation}
Also $\Wc(\beta\Jc(0))=\tfrac{1}{2\beta}W_3(\Jh(0))=\beta_0/\beta=t$ with $t\equiv T/T_0$.
Then \eqref{eq:constraint} reads $t+M_c^2=1$, giving
\begin{equation}
  \boxed{\ \Hc_c(T)=4\Jc(0)\left(1-\frac{T}{T_0}\right)^{1/2}\ }
  \label{eq:boundary}
\end{equation}
i.e.\ $H_c=4\Jh(0)\sqrt{1-t}$, which is (K2.39) at $\lambda=1$. In the control variables,
using $t=\beta_0/\beta$ and $\beta_0\Jc(0)=W_{\rm sc}$, and writing $u\equiv\beta\Jc(0)$,
\begin{equation}
  h_c(\beta)=4u\left(1-\frac{W_{\rm sc}}{u}\right)^{1/2}.
  \label{eq:boundary_bh}
\end{equation}

\paragraph{Topology.} The ordered region is $|h|<h_c(\beta)$ for $\beta>\beta_0$: a
semi-infinite tongue, symmetric in $h\to-h$, attached to the $\beta$ axis at the single
point $u=W_{\rm sc}$ and opening asymptotically linearly, $h_c\simeq4u$. This is
topologically identical to [BR, Fig.~1]: a path from inside to outside must cross the
critical line, since going around it requires $\beta\to\infty$. The motivating claim of
the manuscript survives the replacement, now resting on an exactly solved $d=3$ model.

\subsection{First derivatives}
Because $\psi$ is a stationary value, first derivatives may be taken at fixed $z,x$
(envelope theorem):
\begin{equation}
  \frac{1}{N}\frac{\partial\psi}{\partial h}
   =\frac{\partial\Phi}{\partial h}\bigg|_z=\frac{h}{2a_0}=M ,
\end{equation}
\begin{equation}
  \frac{1}{N}\frac{\partial\psi}{\partial\beta}
   =\frac{\partial\Phi}{\partial\beta}\bigg|_z
   =-\int\!\dk\,\frac{\Jc(\bk)}{a_{\bk}}-\frac{h^2\Jc(0)}{4a_0^2}+\Jc(0)x^2
   =-\left[\int\!\dk\,\frac{\Jc(\bk)}{a_{\bk}}+\Jc(0)M^2-\Jc(0)x^2\right]
   =-\frac{\avg{E_{\rm int}}}{N} \ \checkmark
\end{equation}
the last equality by direct computation of $\avg{E_{\rm int}}$ from \eqref{eq:diag}, whose
Gaussian average is $\avg{|\bm S_{\bk}|^2}=1/a_{\bk}$ plus the uniform condensate $NM^2$ at
$\bk=0$ and the staggered condensate $Nx^2$ at $\bk=\bQ$. The term $-\Jc(0)x^2$ is the
energy of the macroscopic antiferromagnetic mode, since $\Jc(\bQ)=-\Jc(0)$; an earlier
version of these notes omitted it. It does not affect \eqref{eq:psiord}, where the $x^2$
coefficient vanishes identically, nor anything in \S3--\S4, which use fluctuations rather
than the mean energy.

\subsection{The ordered phase in closed form}
\label{sec:ordered}
Here \eqref{eq:sticking} gives $z=\beta\Jc(0)$, so the $x^2$ term in \eqref{eq:Phi} has
\emph{zero coefficient} and the order parameter drops out of $\psi$ entirely. With
$a_{\bk}=\beta[\Jc(0)+\Jc(\bk)]$ and $a_0=2\beta\Jc(0)$,
\begin{equation}
  \boxed{\
  \frac{\psi}{N}\bigg|_{\rm ord}
   =\beta\Jc(0)-\ln\beta+\frac{h^2}{8\beta\Jc(0)}+\mathrm{const}\ }
  \label{eq:psiord}
\end{equation}
the constant being $-\int\!\dk\,\ln[\Jc(0)+\Jc(\bk)]+\ln\pi$, independent of $(\beta,h)$.
Elementary and analytic in $(\beta,h)$ throughout the ordered region. Every static quantity
follows by inspection:
\begin{equation}
  M=\frac{h}{4\beta\Jc(0)},\qquad
  \frac{1}{N}\frac{\partial^2\psi}{\partial h^2}=\frac{1}{4\beta\Jc(0)},\qquad
  \frac{1}{N}\frac{\partial^2\psi}{\partial\beta\partial h}=-\frac{h}{4\beta^2\Jc(0)},\qquad
  \frac{1}{N}\frac{\partial^2\psi}{\partial\beta^2}=\frac{1}{\beta^2}+\frac{h^2}{4\beta^3\Jc(0)} .
  \label{eq:ordhess}
\end{equation}
All finite. This is the promised demonstration that the working rule of
\S\ref{sec:ensembles} costs almost nothing: on the ordered side one never needs the
implicit machinery at all.

\subsection{Second derivatives and the constraint response}
On the disordered side $z$ moves with $\lambda$ and the envelope theorem no longer applies.
Differentiating $\partial\Phi/\partial z=0$ gives
$\partial z/\partial\lambda^\nu=-(\partial^2\Phi/\partial z\partial\lambda^\nu)/(\partial^2\Phi/\partial z^2)$,
hence
\begin{equation}
  \boxed{\
  \frac{1}{N}\frac{\partial^2\psi}{\partial\lambda^\mu\partial\lambda^\nu}
  =\frac{\partial^2\Phi}{\partial\lambda^\mu\partial\lambda^\nu}\bigg|_z
   -\frac{\dfrac{\partial^2\Phi}{\partial\lambda^\mu\partial z}\;
          \dfrac{\partial^2\Phi}{\partial\lambda^\nu\partial z}}
         {\dfrac{\partial^2\Phi}{\partial z^2}}\ }
  \label{eq:hessian}
\end{equation}
which is (K2.38) in general form. The ingredients, from \eqref{eq:Phi}, with
$I_n\equiv\int\!\dk\,a_{\bk}^{-n}$ and $P\equiv\Jc(0)$:
\begin{align}
  \frac{\partial^2\Phi}{\partial z^2}&=I_2+\frac{h^2}{2a_0^3},
  &\frac{\partial^2\Phi}{\partial\beta\partial z}&=\int\!\dk\,\frac{\Jc(\bk)}{a_{\bk}^2}+\frac{h^2P}{2a_0^3},
  &\frac{\partial^2\Phi}{\partial\beta^2}\bigg|_z&=\int\!\dk\,\frac{\Jc(\bk)^2}{a_{\bk}^2}+\frac{h^2P^2}{2a_0^3},\nonumber\\
  \frac{\partial^2\Phi}{\partial h\partial z}&=-\frac{h}{2a_0^2},
  &\frac{\partial^2\Phi}{\partial h\partial\beta}&=-\frac{hP}{2a_0^2},
  &\frac{\partial^2\Phi}{\partial h^2}\bigg|_z&=\frac{1}{2a_0}.
\end{align}

\subsection{The cancellation, explicitly}
\label{sec:cancel}
As $r\to0$ the divergences all come from $\bk$ near $\bQ$. By \eqref{eq:smallq},
\begin{equation}
  I_2=\int\!\dk\,\frac{1}{(r+\tfrac{\beta\Jc}{2}q^2)^2}+\text{reg}
   \;\sim\; \frac{D}{\sqrt r}\to\infty ,
  \qquad
  I_1\to\text{finite},
\end{equation}
(the second because $\int d^3q/q^2$ converges in $d=3$). The trick is to eliminate
$\Jc(\bk)$ in favour of $a_{\bk}$: since $a_{\bk}=r+\beta[\Jc(0)+\Jc(\bk)]$,
\begin{equation}
  \Jc(\bk)=-P+\frac{a_{\bk}-r}{\beta}
  \quad\Longrightarrow\quad
  \int\!\dk\,\frac{\Jc}{a^2}=-PI_2+\frac{A}{\beta},
  \qquad
  \int\!\dk\,\frac{\Jc^2}{a^2}=P^2I_2-\frac{2PA}{\beta}+\frac{C}{\beta^2},
\end{equation}
with $A\equiv I_1-rI_2$ and $C\equiv I_0-2rI_1+r^2I_2$, $I_0=1$. Writing
$R\equiv h^2/2a_0^3$ (regular, since $a_0=r+2\beta P$ stays finite), the three ingredients
become
\begin{equation}
  \frac{\partial^2\Phi}{\partial z^2}=I_2+R,
  \qquad
  \frac{\partial^2\Phi}{\partial\beta\partial z}=-P\big(I_2+R\big)+B,
  \qquad
  \frac{\partial^2\Phi}{\partial\beta^2}\bigg|_z=P^2\big(I_2+R\big)-\frac{2PA}{\beta}+\frac{C}{\beta^2},
\end{equation}
where $B\equiv A/\beta+2PR$ is regular. Substituting into \eqref{eq:hessian}, writing
$D_z\equiv I_2+R$, the $P^2D_z$ terms cancel identically:
\begin{align}
  \frac{1}{N}\frac{\partial^2\psi}{\partial\beta^2}
  &=P^2D_z-\frac{2PA}{\beta}+\frac{C}{\beta^2}-\frac{(-PD_z+B)^2}{D_z}\nonumber\\
  &=P^2D_z-\frac{2PA}{\beta}+\frac{C}{\beta^2}-P^2D_z+2PB-\frac{B^2}{D_z}
   \;=\;\frac{C}{\beta^2}+4P^2R-\frac{B^2}{D_z} .
  \label{eq:cancelled}
\end{align}
\emph{The leading divergence cancels exactly}, because near $\bQ$ the coupling
$\Jc(\bk)\to-P$ is constant, so the singular part of $\partial^2\Phi/\partial\beta\partial z$
is a fixed multiple of that of $\partial^2\Phi/\partial z^2$. Taking $r\to0$ in
\eqref{eq:cancelled}: $C\to I_0=1$, $B^2/D_z\to0$, and $4P^2R=h^2/4\beta^3P$, so
\begin{equation}
  \frac{1}{N}\frac{\partial^2\psi}{\partial\beta^2}\ \longrightarrow\
  \frac{1}{\beta^2}+\frac{h^2}{4\beta^3\Jc(0)} ,
\end{equation}
which is exactly \eqref{eq:ordhess}. The same steps give
$\partial^2\psi/\partial h^2\to1/4\beta\Jc(0)$ and
$\partial^2\psi/\partial\beta\partial h\to-h/4\beta^2\Jc(0)$, again matching
\eqref{eq:ordhess}. \checkmark

\paragraph{Two things this buys.}
\emph{(i)} The Hessian is finite \emph{and continuous} across the critical line --- computed
independently from the two sides. \emph{(ii)} The leading singular correction is
$-B^2/D_z\sim-B^2\sqrt r/D$. Since $r\sim\xi^{-2}\sim\epsilon^{2\nu}=\epsilon^2$, this is
$\sim\epsilon^{1}=\epsilon^{-\alpha}$ with $\alpha=-1$: the expected specific-heat
singularity, obtained here from the explicit cancellation rather than assumed.

\subsection{Rank, and the contrast with the mean-field example}
From \eqref{eq:ordhess},
\begin{equation}
  \det\!\left[\frac{1}{N}\frac{\partial^2\psi}{\partial\lambda\partial\lambda}\right]
  =\left(\frac{1}{\beta^2}+\frac{h^2}{4\beta^3P}\right)\frac{1}{4\beta P}
   -\frac{h^2}{16\beta^4P^2}
  =\frac{1}{4\beta^3\Jc(0)}\;>\;0 .
\end{equation}
Full rank. This is a real structural difference from [BR,\S e], where the disordered-phase
metric is rank one with $\det g=0$ and the optimum collapses onto lines of constant $m$.
There are no null directions here, so the four classes of minimal paths must be re-derived
in \S7, and the ``detour to $\beta=0$'' mechanism probably does not survive.

\subsection{Critical exponents and what they imply}
\label{sec:crit}
The $d=3$ spherical model has, exactly,
\begin{equation}
  \alpha=-1,\quad \beta_{\rm op}=\tfrac12,\quad \gamma=2,\quad \delta=5,\quad
  \nu=1,\quad \eta=0 ,
\end{equation}
($\nu=1/(d-2)$, and $\nu=1$ is confirmed independently by $\xi^2=\beta\Jc/2r$ above), and
with Model~A dynamics $z_{\rm dyn}=2$ exactly. In the notation of [BR, Eq.~(5)],
\begin{equation}
  b=d\nu=2-\alpha=3,\qquad
  \Delta=\beta_{\rm op}\delta=\tfrac52,\qquad
  p=\frac{b}{\Delta}=1+\frac1\delta=\frac65,\qquad
  \frac{z_{\rm dyn}}{d}=\frac23 .
\end{equation}
Both convergence conditions of [BR] hold, exactly rather than numerically:
\begin{equation}
  b\Big(1-\frac{z_{\rm dyn}}{d}\Big)=1>0,
  \qquad
  p\Big(1-\frac{z_{\rm dyn}}{d}\Big)=\frac25>0 .
\end{equation}
\FLAG{Corrected --- $p$ does not describe this control manifold.} An earlier version claimed
$\kappa=pz_{\rm dyn}/d=4/5<1$ ``in the field direction'', offered as a sharper illustration
than the current draft has. That is wrong. The Widom variable $h/t^\Delta$ is built from
the field \emph{conjugate to the order parameter}, which here is the transverse staggered
field $\Hc'_\perp$ --- and $\Hc'_\perp=0$ throughout this control manifold. With
$\Hc'_\perp=0$ we sit on the symmetry axis: at a generic point of the critical line,
variations of $\beta$ and of the \emph{uniform} field $h$ merely have different projections
onto the same thermal scaling field normal to the line, and the tangential variation is
non-singular at leading order. Only $b$ and the condition $b(1-z_{\rm dyn}/d)=1>0$ are
relevant here; $p=6/5$ remains correct universality-class data but describes an approach
this slice cannot realise.

This is not merely a lost selling point --- it is the scaling-language explanation of the
finite metric found in \S4: the control parameters never couple to the ordering field.

\paragraph{The key structural point.} Sections \ref{sec:ordered} and \ref{sec:cancel}
establish by explicit computation that
\begin{equation}
  \boxed{\ \text{the static Hessian does \emph{not} diverge at the transition,}\ }
\end{equation}
consistent with $\alpha=-1<0$ and with [K]'s statement (below K2.60) that $\chi$, $\chi'$
and $C_v$ remain finite and continuous across the boundary for $H\neq0$. Together with the
decoupling \eqref{eq:decoupling}, the entire divergence of $g_{\mu\nu}$ must come from
$T_{\mu\nu}$. This makes the decomposition question decisive:
\begin{align}
  (\text{full Hessian})\times\tau
   &\sim O(1)\cdot\epsilon^{-z_{\rm dyn}\nu}=\epsilon^{-2}
   &&\Rightarrow\ \sqrt g\sim\epsilon^{-1}
   &&\Rightarrow\ \mathcal L\ \text{log-divergent},\\
  (\text{singular})\times\tau+(\text{regular})\times\tau_{\rm micro}
   &\sim \epsilon^{-\alpha}\epsilon^{-z_{\rm dyn}\nu}+O(1)=\epsilon^{-1}
   &&\Rightarrow\ \sqrt g\sim\epsilon^{-1/2}
   &&\Rightarrow\ \mathcal L\ \text{finite}.
\end{align}
The second is correct --- the analytic background relaxes microscopically and cannot be
multiplied by the critical $\tau$ --- and agrees with [BR, Eq.~(11)], which gives
$g_{tt}\dot t^2\propto s^{b(1-z_{\rm dyn}/d)-2}=s^{-1}$ and
$\dot{\mathcal L}\propto s^{-1/2}$. Here the split is not a modelling assumption: the
decomposition in \eqref{eq:cancelled} exhibits the regular piece ($C/\beta^2+4P^2R$) and
the singular piece ($-B^2/D_z$) separately. \S4 must attach the correct relaxation time to
each.

Note $g\sim\epsilon^{-1}$ is the same power as the mean-field $g\propto\Delta^{-1}$ of
[BR,\S e], so the replacement preserves the qualitative picture while upgrading its status.

\paragraph{Proposed new row for Table~I of [BR].}
\begin{center}
\begin{tabular}{lccc}
\hline
 & $b$ & $p$ & $z_{\rm dyn}/d$ \\
\hline
3D spherical (Model A) & $3$ & $6/5$ & $2/3$ \\
\hline
\end{tabular}
\end{center}
All three entries exact.



\medskip
\noindent
\FLAG{Status after referee comments.} The results below are \emph{not} an exact
derivation of the constrained metric, and should be read as a controlled large-$N$
calculation with two known gaps: (i) only the total constraint $Q_+=N$ is imposed, whereas
strict dynamics requires $Q_-=0$ as well (\S1.5 of the \S0--\S2 notes), and (ii) the
constraint is imposed through a static conditioning formula rather than by solving the
projected dynamics. The energy-weight identity \eqref{eq:identity} is exact and is the
physically important content. See \S\ref{sec:whatisproved}.

\medskip
\noindent
\textbf{Summary of what \S3--\S4 suggest.} The soft modes decouple from both control
parameters, exactly, as a consequence of the spherical constraint. Consequently the
dissipation metric is \emph{finite and continuous across the phase transition} --- the
diverging relaxation time is multiplied by a susceptibility that vanishes at precisely the
compensating rate. In the ordered phase the metric is not merely finite but \emph{flat}:
the ordered region is a Euclidean half-disk of radius $2$ in the coordinates
$(2\sqrt{T/T_0},\,2M)$, with the critical line as its boundary arc.

The flatness, however, is specific to one choice of dissipation functional (\S4.1), and
the crossing statement is an argument rather than a theorem (\S\ref{sec:crossing}).

%===============================================================
\section{Model A dynamics}
%===============================================================

\subsection{Why non-conserved}
The protocol varies $h=\beta\Hc$, and $\partial\psi/\partial h=\avg{M_{\rm tot}}$: driving
the field \emph{must} change the uniform magnetisation. A dynamics conserving
$M_{\rm tot}$ is therefore incompatible with the control problem, and Model~B is excluded
on physical grounds rather than by convention. This also matches the Glauber (single-spin,
non-conserved) choice made in [BR, End Matter].

The order parameter here is the staggered transverse magnetisation, which is \emph{not}
conserved by the uniform field either, so Model~A is the appropriate choice for it as well.
With $\eta=0$ the dynamic exponent is then $z_{\rm dyn}=2-\eta=2$ exactly, confirmed
directly below.

\subsection{Langevin equations at fixed spherical field}
Write the dimensionless effective Hamiltonian $\beta\Hc_{\rm eff}=\sum_{\bk}a_{\bk}|\bm S_{\bk}|^2-h\sqrt N S^z_0$.
Model~A reads
\begin{equation}
  \partial_t \bm S_{\bk}
   = -\Gamma\,\frac{\partial(\beta\Hc_{\rm eff})}{\partial \bm S_{\bk}^{*}}+\bm\eta_{\bk},
  \qquad
  \avg{\eta^{a}_{\bk}(t)\,\eta^{b*}_{\bk'}(t')}=2\Gamma\,\delta^{ab}\delta_{\bk\bk'}\delta(t-t'),
\end{equation}
with $\Gamma$ the microscopic rate. Explicitly,
\begin{equation}
  \boxed{\ \partial_t\bm S_{\bk}=-2\Gamma a_{\bk}\bm S_{\bk}
   +\Gamma h\sqrt N\,\delta_{\bk,0}\,\hat e_z+\bm\eta_{\bk}\ }
  \label{eq:langevin}
\end{equation}
so each mode is an independent Ornstein--Uhlenbeck process with
\begin{equation}
  \text{relaxation rate}\quad \Gamma_{\bk}=2\Gamma a_{\bk},
  \qquad
  \text{stationary variance}\quad \avg{|\delta S^a_{\bk}|^2}=\frac{1}{2a_{\bk}} ,
  \label{eq:ourates}
\end{equation}
the variance being the equilibrium one, as required by detailed balance
($2\Gamma/(2\cdot 2\Gamma a_{\bk})=1/2a_{\bk}$). \checkmark

\paragraph{Note on factors of two.} \eqref{eq:langevin} carries $2\Gamma a_{\bk}$, not
$\Gamma a_{\bk}$. Since $\Gamma$ is an arbitrary microscopic rate this is pure convention,
but it must be used consistently: bilinear observables then relax at $2\Gamma_{\bk}=4\Gamma a_{\bk}$,
and every entry of the metric in \S4 inherits a $1/\Gamma$ with a definite numerical
coefficient.

\subsection{The dynamic critical exponent}
Near the zone corner, using $a_{\bQ-\bq}=r+\tfrac{\beta\Jc}{2}q^2$ from \S1.5,
\begin{equation}
  \Gamma_{\bQ-\bq}=2\Gamma\left(r+\tfrac{\beta\Jc}{2}q^2\right),
  \qquad
  \tau_{\bQ}=\frac{1}{2\Gamma r} .
\end{equation}
With $r\propto\xi^{-2}$ (\S2.3) this gives $\tau\propto\xi^{2}$, i.e.
\begin{equation}
  \boxed{\ z_{\rm dyn}=2\ }
\end{equation}
and since $\nu=1$, $\tau\sim\epsilon^{-z_{\rm dyn}\nu}=\epsilon^{-2}$. The relaxation time
genuinely diverges at the transition; the content of \S4 is that this does not, by itself,
make the metric diverge.

\subsection{The constraint must be imposed instantaneously}
\label{sec:instantaneous}
Equation \eqref{eq:langevin} at fixed $z$ equilibrates to the fixed-$z$ Gaussian measure,
which by \S2.3 is \emph{not} the physical one: it has superextensive fluctuations of
$Q\equiv\sum_{\bk}|\bm S_{\bk}|^2$ and gives the wrong (divergent) response functions. The
dynamics must therefore enforce the constraint on the configuration itself,
\begin{equation}
  Q(t)=\sum_{\bk}|\bm S_{\bk}(t)|^2=N \quad\text{for all }t ,
  \label{eq:hardconstraint}
\end{equation}
by promoting $z$ to a fluctuating Lagrange multiplier $\mu(t)$ in \eqref{eq:langevin}.

\FLAG{Revision to \S1.3.} This qualifies the earlier statement that the mean spherical
ensemble is ``good enough''. It is, for statics --- the ensembles are thermodynamically
equivalent and \S2 is unaffected. But the metric is built from fluctuations, and the
dynamics that generates the correct fluctuations is the hard-constraint one. Imposing
\eqref{eq:hardconstraint} only on average would leave the fixed-$z$ spectrum in place and
produce a metric that diverges over the whole ordered region.

Differentiating \eqref{eq:hardconstraint} and using \eqref{eq:langevin} gives $\mu(t)$
explicitly,
\begin{equation}
  \mu(t)=z+\delta\mu(t),
  \qquad
  \delta\mu = \frac{1}{2\Gamma N}\Big[\textstyle\sum_{\bk}\xi_{\bk}\Big]
              -\frac{1}{N}\sum_{\bk}a_{\bk}\,\delta|\bm S_{\bk}|^2 ,
  \label{eq:mu}
\end{equation}
where $\xi_{\bk}$ collects the noise contributions to $\dot Q$.

\FLAG{Three caveats.} (i) \eqref{eq:mu} omits the source-dependent term
$\propto h\,\delta S^z_0/\sqrt N$ arising from the drive at $\bk=0$.
(ii) $\delta\mu$ is \emph{not} $O(1/N)$: both the extensive quadratic fluctuation and the
noise sum are $O(\sqrt N)$, so $\delta\mu=O(N^{-1/2})$ as a fluctuating quantity. A
multiplier of that size still acts at order one on extensive observables, so smallness
alone does not justify neglecting it. (iii) Differentiating $Q(t)$ for white noise requires
committing to It\^o or Stratonovich; the quadratic-variation term is missing above. What
does survive is that the systematic part of $\delta\mu$ is proportional to
$\delta E_{\rm int}$ itself, and acts at the non-critical rate $O(\Gamma)$.

\subsection{The key identity: the soft modes carry no energy weight}
\label{sec:identity}
This is the central structural fact of these notes, and the dynamical counterpart of the
static cancellation of \S2.9.

From $a_{\bk}=z+\beta\Jc(\bk)$ we have $\Jc(\bk)=(a_{\bk}-z)/\beta$ identically, so
\begin{equation}
  E_{\rm int}=\sum_{\bk}\Jc(\bk)|\bm S_{\bk}|^2
   =\frac{1}{\beta}\sum_{\bk}a_{\bk}|\bm S_{\bk}|^2-\frac{z}{\beta}\underbrace{\sum_{\bk}|\bm S_{\bk}|^2}_{=\,N}
   =\frac{\Ac}{\beta}-\frac{zN}{\beta},
  \qquad
  \Ac\equiv\sum_{\bk}a_{\bk}|\bm S_{\bk}|^2 .
\end{equation}
The second term is a constant \emph{on the constraint surface}. Hence, exactly,
\begin{equation}
  \boxed{\ \delta E_{\rm int}=\frac{1}{\beta}\,\delta\Ac
   =\frac{1}{\beta}\sum_{\bk}a_{\bk}\,\delta|\bm S_{\bk}|^2\ }
  \label{eq:identity}
\end{equation}

The weight attached to mode $\bk$ is $a_{\bk}$, which \emph{vanishes at the soft mode}.
Precisely: the zero-momentum critical mode $\bk=\bQ$ has exactly zero energy weight, and
nearby critical modes enter only through the gradient weight $a_{\bQ-\bq}\simeq r+cq^2$.
They are strongly suppressed, not absent.
Physically: a long-wavelength fluctuation of the staggered order parameter costs no energy,
so it cannot appear in energy fluctuations. The spherical constraint enforces this exactly,
because near $\bQ$ the coupling $\Jc(\bk)\to-P$ is constant and a constant weight is
precisely what the constraint removes.

\paragraph{What this does to the divergence.} Naively, from $E_{\rm int}=\sum\Jc|\bm S|^2$,
\begin{equation}
  \int_0^\infty\!\!\avg{\delta E(t)\delta E(0)}dt
  \ \sim\ \sum_{\bk}\frac{\Jc(\bk)^2}{a_{\bk}^2}\cdot\frac{1}{a_{\bk}}
  \ \sim\ \int\!\frac{d^3q}{q^6}\qquad\text{(divergent)},
\end{equation}
which is the obstruction flagged at the end of the \S0--\S2 notes. Using
\eqref{eq:identity} instead,
\begin{equation}
  \int_0^\infty\!\!\avg{\delta E(t)\delta E(0)}dt
  \ \sim\ \sum_{\bk}\frac{a_{\bk}^2}{a_{\bk}^2}\cdot\frac{1}{a_{\bk}}
  \ \sim\ \int\!\frac{d^3q}{q^2}\qquad\text{(convergent in $d=3$)} .
\end{equation}
Four powers of $q$ are gained. The obstruction is removed, not by a delicate cancellation
of large terms, but by the observable simply not containing the dangerous modes.

\subsection{Why the residual constraint corrections are negligible}
\label{sec:residual}
Having rewritten the observable, one may compute with the \emph{fixed-$z$} Gaussian
dynamics \eqref{eq:langevin}, provided the residual effect of the constraint is small. The
tool used below is the Gaussian conditioning formula
\begin{equation}
  \text{[large-$N$, not an identity]}\qquad\qquad\qquad\qquad\qquad\qquad\qquad\nonumber
\end{equation}
\FLAG{Not exact.} $Q=\sum_{\bk}|\bm S_{\bk}|^2$ is \emph{quadratic} in the Gaussian
variables, so $A$, $B$, $Q$ are not jointly Gaussian and \eqref{eq:conditional} is a
large-$N$ regression, not an exact conditioning identity. Its justification here is
a posteriori: it reproduces the exact static Schur complement of \S2.8 (see \S4.3), which
was derived independently by implicit differentiation of the stationary $\psi$. That is
good evidence it is correct to leading order in $N$, but it is not a proof, and inserting a
\emph{static} projector into a time-integrated correlation function is a further
uncontrolled step.
\begin{equation}
  \mathrm{Cov}(A,B\,|\,Q)=\mathrm{Cov}(A,B)-\frac{\mathrm{Cov}(A,Q)\,\mathrm{Cov}(B,Q)}{\mathrm{Var}(Q)} .
  \label{eq:conditional}
\end{equation}
With $I_p\equiv\int\!\dk\,a_{\bk}^{-p}$ (so $I_1=\Wc$, $I_0=1$), and counting $nN$ real
degrees of freedom each contributing $\mathrm{Var}(u^2)=1/2a^2$,
\begin{equation}
  \frac{\mathrm{Var}(\Ac)}{N}=\frac{n}{2},
  \qquad
  \frac{\mathrm{Cov}(\Ac,Q)}{N}=\frac{n}{2}I_1,
  \qquad
  \frac{\mathrm{Var}(Q)}{N}=\frac{n}{2}I_2 ,
\end{equation}
so that
\begin{equation}
  \frac{\mathrm{Cov}(\Ac,\Ac|Q)}{N}=\frac{n}{2}\left[1-\frac{I_1^2}{I_2}\right].
  \label{eq:projectedvar}
\end{equation}
In the ordered phase $r=0$ and $I_2=\int\!\dk\,a_{\bk}^{-2}\sim\int d^3q/q^4$ diverges, so
the correction vanishes identically. Approaching from the disordered side
$I_2\simeq D/\sqrt r\to\infty$ and the correction is $O(\sqrt r)$.

The same holds for the time-integrated quantities, with $I_2\to I_2$ and $I_3$ replacing the
static moments: the correction scales as $I_2^2/I_3\sim r^{-1}/r^{-3/2}=\sqrt r\to0$.

\CHECK{This last estimate treats the constrained dynamics through the static projector
\eqref{eq:conditional}. The exact constrained process has the drift feedback of
\eqref{eq:mu}, whose systematic part is $\propto\delta E$ itself and acts at the
\emph{non-critical} rate $O(\Gamma)$; it therefore renormalises the fast sector but cannot
generate a critical divergence. Worth confirming by diagonalising the constrained linear
dynamics explicitly before this is committed to the manuscript.}

\subsection{Correlation functions}
\FLAG{Gap.} What is separated below is the \emph{uniform} mean $\avg{S^z_0}=\sqrt N M$,
which is not the antiferromagnetic condensate. The latter is
$S^x_{\bQ}=\sqrt N\,x+\delta S^x_{\bQ}$, and in the ordered phase $a_{\bQ}=0$, so the
fixed-$z$ Ornstein--Uhlenbeck process \eqref{eq:ourates} has no restoring force there and
no stationary distribution: \eqref{eq:onepart}--\eqref{eq:twopart} cannot be applied to a
sum that includes $\bQ$ without separating it first. By \eqref{eq:identity} the $\bQ$-mode
carries zero energy weight and so should contribute nothing directly to $\delta E$, and it
does not couple to $M$ at linear order --- but its role in enforcing the two constraints,
and its effect on the dynamics of the remaining modes, has to be derived rather than
assumed. This is the main missing piece.

Separate the uniform mean. With $\avg{S^z_0}=m_0=\sqrt N M$, write
$\bm S_0=(S^x_0,\;m_0+v)$ with $v\equiv\delta S^z_0$. Then from \eqref{eq:identity}
\begin{equation}
  \delta\Ac = \underbrace{2a_0m_0\,v}_{\text{one-particle}}
   \;+\;\underbrace{\Sigma}_{\text{two-particle}},
  \qquad
  \Sigma\equiv\sum_{\bk}a_{\bk}\,\delta|\bm S_{\bk}|^2\Big|_{\text{fluct.}} ,
\end{equation}
and $\delta X_2=\delta M_{\rm tot}=\sqrt N\,v$. The two sectors are uncorrelated (odd
versus even in the Gaussian variables), which is what makes the metric assemble so simply.
Their correlators follow from \eqref{eq:ourates}:
\begin{align}
  \avg{v(t)v(0)}&=\frac{1}{2a_0}\,e^{-2\Gamma a_0 t},
  &&\int_0^\infty\!\!=\frac{1}{4\Gamma a_0^2},
  \label{eq:onepart}\\[2pt]
  \avg{\Sigma(t)\Sigma(0)}&=\frac{n}{2}\sum_{\bk}e^{-4\Gamma a_{\bk}t},
  &&\int_0^\infty\!\!=\frac{n}{8\Gamma}\sum_{\bk}\frac{1}{a_{\bk}}=\frac{nN}{8\Gamma}\,\Wc(z),
  \label{eq:twopart}
\end{align}
the bilinears relaxing at twice the single-mode rate, and their weight $a_{\bk}^2$
cancelling the variance $1/2a_{\bk}^2$ mode by mode --- the identity \eqref{eq:identity}
at work. Note the appearance of $\Wc(z)$, which the constraint \eqref{eq:constraint} of
\S2.3 fixes in closed form:
\begin{equation}
  \Wc(z)=1-M^2-x^2 .
  \label{eq:Wclosed}
\end{equation}

%===============================================================
\section{The dissipation metric}
%===============================================================

\subsection{Which functional is being minimised}
\label{sec:objective}
This is not a matter of convention and must be settled before any geometric statement is
made, because the two natural objectives differ by a non-constant conformal factor.

\paragraph{(a) Excess entropy production.} In natural Massieu coordinates
$\partial_\mu\psi=\avg{X_\mu}$, $\partial_\mu\partial_\nu\psi=\avg{\delta X_\mu\delta X_\nu}$
(the Fisher information of the exponential family), and the excess entropy production of a
slow protocol is $\Delta S_{\rm ex}=\int g^{S}_{\mu\nu}\dot\lambda^\mu\dot\lambda^\nu dt$ with
\begin{equation}
  \boxed{\ g^{S}_{\mu\nu}=\int_0^\infty \avg{\delta X_\mu(t)\,\delta X_\nu(0)}\,dt\ }
  \label{eq:metricdef}
\end{equation}
and no prefactor.

\paragraph{(b) Dissipated availability, as in [BR, Eq.~(1)].} There
$\zeta_{ij}=\beta\int\avg{\delta X_i\delta X_j}$ with $\lambda$ energy-like. Transforming
to our coordinates is a tensor operation, e.g.\ at fixed $\beta$,
\begin{equation}
  \zeta_{hh}=\zeta_{HH}\Big(\frac{\partial \Hc}{\partial h}\Big)^{2}
   =\beta\!\int\!\avg{\delta M\delta M}\cdot\frac{1}{\beta^{2}}
   =\frac{1}{\beta}\int\!\avg{\delta M\delta M},
\end{equation}
so
\begin{equation}
  g^{A}_{\mu\nu}=\frac{1}{\beta}\,g^{S}_{\mu\nu},
\end{equation}
consistent with $A\simeq T\,\Delta S_{\rm ex}$.

\FLAG{Correction.} An earlier version of these notes asserted that dropping the $\beta$ was
merely a change of coordinates. It is not: a conformal factor cannot be removed by
reparametrisation, and the two objectives give genuinely different geometries. Everything
below uses (a). Section \ref{sec:flat} records what (b) gives.

All quantities below are per site.

\subsection{General result}
Assembling \eqref{eq:onepart}--\eqref{eq:twopart} with $\delta X_1=-\delta\Ac/\beta$ and
$\delta X_2=\sqrt N v$:
\begin{align}
  g_{hh}&=\frac{1}{4\Gamma a_0^{2}},
  \label{eq:ghh}\\
  g_{\beta h}&=-\frac{M}{2\Gamma\beta a_0},
  \label{eq:gbh}\\
  g_{\beta\beta}&=\frac{1}{\Gamma\beta^{2}}\left[M^{2}+\frac{\Wc(z)}{4}\right]
   \;=\;\frac{1}{4\Gamma\beta^{2}}\Big[\,4M^{2}+1-M^{2}-x^{2}\Big],
  \label{eq:gbb}
\end{align}
using \eqref{eq:Wclosed} in the last step. Every entry is manifestly finite wherever
$a_0>0$, which holds throughout the phase diagram.

\subsection{Static check}
Setting $t=0$ in the correlators must reproduce the Hessian of \S2. The equal-time
covariances are
\begin{equation}
  \frac{\avg{\delta X_1^2}}{N}=\frac{1}{\beta^2}\Big[2a_0M^2+\tfrac n2\Big],
  \qquad
  \frac{\avg{\delta X_1\delta X_2}}{N}=-\frac{M}{\beta},
  \qquad
  \frac{\avg{\delta X_2^2}}{N}=\frac{1}{2a_0}.
\end{equation}
In the ordered phase $a_0=2\beta P$ and $M=h/4\beta P$, so with $n=2$
\begin{equation}
  \frac{\avg{\delta X_1^2}}{N}=\frac{1}{\beta^2}\Big[\tfrac{h^2}{4\beta P}+1\Big]
   =\frac{1}{\beta^2}+\frac{h^2}{4\beta^3\Jc(0)},
  \qquad
  -\frac{M}{\beta}=-\frac{h}{4\beta^2\Jc(0)},
  \qquad
  \frac{1}{2a_0}=\frac{1}{4\beta\Jc(0)},
\end{equation}
which are exactly the three entries of Eq.~(2.20) of the \S0--\S2 notes. \checkmark

On the disordered side the projector \eqref{eq:conditional} must be included, and it
reproduces \S2.8 identically: e.g.
\begin{equation}
  \frac{1}{N}\frac{\partial^2\psi}{\partial h^2}
  =\frac{1}{2a_0}-\frac{M^2/a_0^2}{I_2+2M^2/a_0} ,
\end{equation}
which is Eq.~(2.21) of those notes once $h=2a_0M$ is substituted. \checkmark\ This is a
strong check that the fluctuation route and the implicit-differentiation route agree.

\subsection{The ordered phase}
Here $z=\beta P$, $a_0=2\beta P$, and $\Wc(z)=\beta_0/\beta=t$ exactly (evaluate
$\Wc$ at the sticking value and use $\bk\to\bQ-\bk$; this is the same integral that defines
$T_0$). Therefore
\begin{equation}
  g_{\beta\beta}=\frac{4M^{2}+t}{4\Gamma\beta^{2}},
  \qquad
  g_{\beta h}=-\frac{M}{4\Gamma\beta^{2}P},
  \qquad
  g_{hh}=\frac{1}{16\Gamma\beta^{2}P^{2}} ,
  \label{eq:ordmetric}
\end{equation}
\begin{equation}
  \det g=\frac{t}{64\,\Gamma^{2}\beta^{4}P^{2}}\;>\;0 .
\end{equation}
Positive definite and full rank everywhere except $T\to0$. Contrast [BR,\S e], where the
disordered-phase metric is rank one with $\det g=0$; there are no null directions here.

\subsection{The ordered phase is flat}
\label{sec:flat}
Completing the square in \eqref{eq:ordmetric},
\begin{equation}
  ds^{2}=\frac{1}{4\Gamma\beta^{2}}
   \left[\,t\,d\beta^{2}+\Big(2M\,d\beta-\frac{dh}{2P}\Big)^{2}\right].
\end{equation}
Since $M=h/4\beta P$ in the ordered phase,
$dM=\beta^{-1}\big(\tfrac{dh}{4P}-M\,d\beta\big)$, so the bracketed second term is exactly
$4\beta^{2}dM^{2}$ and
\begin{equation}
  ds^{2}=\frac{t\,d\beta^{2}}{4\Gamma\beta^{2}}+\frac{dM^{2}}{\Gamma} .
\end{equation}
Trading $\beta$ for $t=\beta_0/\beta$ gives $t\,d\beta^2/\beta^2=dt^2/t$, hence
\begin{equation}
  ds^{2}=\frac{1}{4\Gamma}\left[\frac{dt^{2}}{t}+4\,dM^{2}\right]
   =\frac{1}{4\Gamma}\Big[dw^{2}+dv^{2}\Big],
  \qquad
  \boxed{\ w\equiv2\sqrt{t}=2\sqrt{T/T_0},\qquad v\equiv 2M\ }
  \label{eq:flat}
\end{equation}

The ordered phase is \textbf{Euclidean}. Moreover the constraint
$t+M^{2}+x^{2}=1$ with $x^{2}>0$ reads
\begin{equation}
  w^{2}+v^{2}<4 ,
\end{equation}
so the ordered region is the half-disk of radius $2$ ($w>0$), the critical line
\eqref{eq:boundary} is its boundary \emph{arc}, and the order parameter is
$x^{2}=1-\tfrac14(w^{2}+v^{2})$: it measures the Euclidean distance from the boundary.
Geodesics are straight lines and thermodynamic length is Euclidean distance divided by
$\sqrt{4\Gamma}$.

No numerical check of the Riemann tensor is needed: the explicit coordinate change to
$(w,v)$ already proves the candidate metric is flat. What is \emph{not} established is that
\eqref{eq:ordmetric} is the correct strictly constrained metric --- see
\S\ref{sec:whatisproved}.

\paragraph{Interpretation of $x^2$.} From $x^{2}=1-\rho^{2}/4$ with
$\rho=\sqrt{w^{2}+v^{2}}$, one has $x^{2}=(2-\rho)(2+\rho)/4$, so $x^2$ is proportional to
the Euclidean distance $2-\rho$ from the boundary arc only asymptotically as $\rho\to2$, not
identically. (Earlier phrasing here was loose.)

\paragraph{With the availability metric (b) instead.} $g^{A}=g^{S}/\beta$ and
$1/\beta=T_0 t/\,\cdot\propto w^{2}$, so
\begin{equation}
  ds^{2}_{A}\ \propto\ w^{2}\big(dw^{2}+dv^{2}\big),
\end{equation}
conformally Euclidean with factor $e^{2\phi}$, $\phi=\ln w$. Since
$K=-e^{-2\phi}\Delta\phi$ and $\Delta\ln w=-1/w^{2}$, this has
\begin{equation}
  K=+\frac{1}{w^{4}}\;>\;0 ,
\end{equation}
positive and non-constant --- \emph{not} flat. The metric degenerates at $w=0$ ($T=0$),
which is physically reasonable for an availability-type objective ($A\simeq T\Delta S_{\rm ex}\to0$),
but it means the elegant Euclidean picture belongs to objective (a) alone. Whichever the
manuscript adopts, it should say so explicitly.

\subsection{The disordered phase}
Here $x=0$, so $\Wc=1-M^{2}$ and
\begin{equation}
  g_{\beta\beta}=\frac{1+3M^{2}}{4\Gamma\beta^{2}},
  \qquad
  g_{\beta h}=-\frac{M}{2\Gamma\beta a_0},
  \qquad
  g_{hh}=\frac{1}{4\Gamma a_0^{2}},
  \qquad
  \det g=\frac{1-M^{2}}{16\,\Gamma^{2}\beta^{2}a_0^{2}} ,
\end{equation}
with $a_0=r+2\beta P$ and $r(\beta,h)$ from the constraint. The same completion of the
square gives
\begin{equation}
  ds^{2}=\frac{1}{4\Gamma}\left[\frac{\Wc}{\beta^{2}}d\beta^{2}
   +\Big(\frac{2M}{\beta}d\beta-\frac{dh}{a_0}\Big)^{2}\right],
\end{equation}
but now $\tfrac{dh}{a_0}-\tfrac{2M}{\beta}d\beta=2\,dM-2M\,d\ln(\beta/a_0)$ and the second
term is no longer a perfect differential, because $a_0/\beta$ is not constant. The
disordered phase is therefore \emph{not} flat in general; it becomes flat only on approach
to the critical line, where $r\to0$ and $a_0/\beta\to2P$.

\subsection{Continuity across the transition}
\label{sec:continuity}
On the critical line, $r=0$ and $x=0$ simultaneously, so $a_0=2\beta P$ and $\Wc=t=1-M_c^{2}$
from both sides. Substituting into \eqref{eq:ghh}--\eqref{eq:gbb}, the ordered-side
expressions \eqref{eq:ordmetric} and the disordered-side expressions coincide:
\begin{equation}
  g_{\beta\beta}\to\frac{1+3M_c^{2}}{4\Gamma\beta^{2}}=\frac{t+4M_c^{2}}{4\Gamma\beta^{2}},
  \qquad
  g_{\beta h}\to-\frac{M_c}{4\Gamma\beta^{2}P},
  \qquad
  g_{hh}\to\frac{1}{16\Gamma\beta^{2}P^{2}} .
\end{equation}
\begin{equation}
  \boxed{\ \text{The metric is finite and continuous across the phase transition.}\ }
\end{equation}
The leading singular behaviour is the $O(\sqrt r)=O(\epsilon)$ correction of
\S\ref{sec:residual}, i.e.\ $\epsilon^{-\alpha}$ with $\alpha=-1$: a vanishing cusp, not a
divergence.

\subsection{Why the divergent relaxation time does not produce a divergent metric}
The relaxation time genuinely diverges, $\tau\sim\epsilon^{-2}$. The reason the metric does
not is the identity \eqref{eq:identity}: the modes whose relaxation diverges enter
$\delta E_{\rm int}$ with weight $a_{\bq}\to0$, and enter $\delta M_{\rm tot}$ not at all
(the field acts only at $\bk=0$, and $a_0$ is non-critical). The critical modes are simply
invisible to both control parameters.

This is the same statement as the symmetry decoupling flagged in \S1.5 of the \S0--\S2
notes, now made quantitative: with $H'=0$ one drives the system through the transition
using fields that couple only to non-critical observables. In scaling language, the
susceptibility that multiplies $\tau$ is the \emph{singular} one, which vanishes as
$\epsilon^{-\alpha}=\epsilon^{1}$; the product $\epsilon^{1}\cdot\epsilon^{-2}=\epsilon^{-1}$
is what [BR, Eq.~(11)] predicts, but that singular piece sits on top of a regular
$O(1)$ background and --- crucially --- the constraint suppresses its coefficient to zero
rather than leaving it finite.

\CHECK{The tension with [BR, Eq.~(11)] is the single most important thing to resolve.
[BR] predicts $g_{tt}\dot t^{2}\propto s^{b(1-z_{\rm dyn}/d)-2}=s^{-1}$, i.e.\ a divergent
metric with integrable $\sqrt g$; the exact calculation here gives a finite metric.
Either the model evades the scaling argument (because the control fields do not couple to
the order parameter, so the ``singular part'' of the metric has vanishing amplitude), or
one of the two is wrong. This must be settled, and if the former, it is worth saying
explicitly in the manuscript --- the scaling criterion would then be sufficient but not
necessary for finite length.}

%===============================================================
\section{What is and is not established}
%===============================================================
\label{sec:whatisproved}

\paragraph{Solid.} The strict-constraint identity \eqref{eq:identity}, and with it the
mechanism: the critical modes have strongly suppressed overlap with both driven
observables, the $\bQ$-mode exactly so. This is the physically valuable content and it
survives every caveat below.

\paragraph{Plausible but not derived.} That the metric is finite and continuous across the
transition. The power counting supports it; the derivation does not yet establish it.

\paragraph{Not established.}
\begin{enumerate}
  \item That \eqref{eq:ghh}--\eqref{eq:gbb} are the exact dissipation metric. They are the
        large-$N$ \emph{unconstrained reference} metric. \S4.3 itself notes that on the
        disordered side the fixed-$z$ equal-time covariance does \emph{not} reproduce the
        Hessian without the projector; the time-integrated version cannot then be exact.
        The physical metric should carry a frequency-dependent low-rank correction --- rank
        two, once $Q_-$ is included. Relatedly, the $O(\sqrt r)$ cusp claimed in \S4.7 lives
        in that omitted correction, not in the displayed expressions, which depend on $r$
        analytically through $a_0$ and $M$ and so give only $O(r)$.
  \item That the ordered metric is exactly flat for the physical dynamics. The coordinate
        change proves \eqref{eq:ordmetric} is flat; it does not prove \eqref{eq:ordmetric}.
        And flatness holds only for objective (a) of \S\ref{sec:objective}.
  \item That the minimal disordered-to-disordered path crosses the ordered phase.
\end{enumerate}

\paragraph{The calculation that would settle it.} Start from a strict, \emph{two}-constraint
projected Model~A process; separate the antiferromagnetic condensate $S^x_{\bQ}$ explicitly;
linearise about it; and compute the zero-frequency response including both fluctuating
spherical multipliers. That determines whether the finite-metric and flat-half-disk picture
is exact, asymptotic, or an artefact of the unconstrained reference.

%===============================================================
\section{Consequences for the manuscript}
%===============================================================
\label{sec:consequences}

The result of \S\ref{sec:continuity} is clean, exact, and pretty. It is also the wrong
result for the paper as currently framed.

\paragraph{What breaks.} The manuscript's central claim is that the metric diverges at a
second-order transition while the thermodynamic length across it can remain finite. In this
model, at $H'=0$, the metric does not diverge at all. The example therefore cannot
illustrate the phenomenon; it illustrates a different (and weaker) one.

\paragraph{What survives --- and is in fact stronger.} The cross-and-return phenomenon of
[BR,\S e] \emph{does} occur here, with the roles of the two phases exchanged, and it can be
proved rather than computed. See \S\ref{sec:crossing}.

\paragraph{The diagnosis.} The loss traces to a single cause: with $H'=0$ neither control
parameter couples to the order parameter. This was flagged as a ``cost'' in \S1.4 of the
\S0--\S2 notes. It costs the divergence, but not the optimal-path result.

\subsection{The optimal path between two disordered points crosses the ordered phase}
\label{sec:crossing}

Set $4\Gamma=1$ for brevity. The half-disk is convex, so two points \emph{in} the ordered
phase are joined by a chord lying entirely inside it: no cross-and-return from that side.
The interesting case is the opposite one.

\paragraph{Geometry of the disordered phase.} With $x=0$ the constraint gives
$\Wc=1-M^{2}$ exactly. Using $(\ln\beta,M)$ as coordinates, with $\theta=z/\beta$ fixed by
$F(\theta)=\beta(1-M^{2})$, $F(\theta)\equiv\int\!\dk\,(\theta+\Jc(\bk))^{-1}$ and
$G\equiv-F'=\beta^{2}I_{2}$,
\begin{equation}
  ds^{2}=(1-M^{2})\,(d\ln\beta)^{2}
        +\Big(2\,dM+\frac{2M\,d\theta}{\theta+P}\Big)^{2},
  \qquad
  d\theta=-\frac{(1-M^{2})\,d\beta-2\beta M\,dM}{G} .
  \label{eq:disordered_line}
\end{equation}
At the critical line $G\to\infty$, the $d\theta$ term switches off, and
\eqref{eq:disordered_line} matches the ordered-phase form with $t\to1-M^{2}$ --- the
continuity of \S\ref{sec:continuity} again.

\paragraph{The claim (an argument, not a theorem).} Let $A=(\beta_1,h_1)$ and
$B=(\beta_1,-h_1)$ with $M_1>0$, both in the disordered phase just outside the arc, and
$\beta_1>\beta_0$ so that the ordered tongue separates them. The minimal path appears to
cross the ordered phase. Two caveats are recorded at the end.

\emph{Through.} In the ordered phase $ds^{2}=dw^{2}+dv^{2}$ with $v=2M$. The chord at
constant $w=w_1$ from $v_1$ to $-v_1$ lies inside the disk, since
$w_1^{2}+v^{2}\le w_1^{2}+v_1^{2}=4$ for $|v|\le|v_1|$. Its length is exactly
\begin{equation}
  L_{\rm through}=2v_1=4M_1
\end{equation}
in the limit $A,B\to$ arc.

\emph{Around.} Remaining in the disordered phase requires $t>1-M^{2}$ pointwise, i.e.
$\beta<\beta_0/(1-M^{2})$. At $M=0$ this forces $\beta<\beta_0$: any disordered path from
$M_1$ to $-M_1$ must pass over the tip of the tongue. There is no route at large $\beta$
either, since $h_c\simeq4\beta P$ grows and $h=0$ stays interior. From
\eqref{eq:disordered_line}, $ds\ge2|dM|$, so
\begin{equation}
  L_{\rm around}\ \ge\ \int 2|dM|\ =\ 4M_1 ,
\end{equation}
with equality only if $d\ln\beta\equiv0$ along the entire path --- which the required
excursion to $\beta<\beta_0$ forbids. Hence
\begin{equation}
  L_{\rm around}>L_{\rm through}
\end{equation}
for endpoints sufficiently close to the critical line.

\CHECK{Two gaps. (i) The bound $ds\ge2|dM|$ is safe when $d\ln\beta=0$, where the $d\theta$
correction multiplies $2dM$ by $1+4\beta M^{2}/[G(\theta+P)]>1$; but along a general path
with $d\beta\neq0$ the $d\theta$ term can carry the opposite sign to $dM$, so the bound is
not airtight and the equality conditions must be obtained from the full quadratic form.
(ii) \eqref{eq:disordered_line} descends from \eqref{eq:gbb}, which is the large-$N$
unconstrained-reference metric, not the established constrained one. The quantitative
excess $2\sqrt{1-M_1^{2}}\ln(\beta_1/\beta_0)$ quoted in an earlier version is not
established by the displayed inequalities and has been withdrawn.}

\paragraph{Robustness to the objective.} The conclusion does not depend on the choice in
\S\ref{sec:objective}. With the availability metric $g^{A}=g^{S}/\beta$, the cost of the
high-temperature detour is $\int\!\sqrt{\Wc/\beta}\,|d\ln\beta|=\int\!\sqrt{\Wc}\,\beta^{-3/2}d\beta$,
which \emph{diverges} as $\beta\to0$: infinite temperature becomes infinitely far, and going
around is even more expensive. This is the opposite of the mean-field antiferromagnet of
[BR,\S e], where $\beta=0$ is the cheapest point of the manifold.

\paragraph{Inverted mechanism.} In [BR,\S e] the crossing is expensive (the metric diverges
on the line) but the far side is cheap, and the content of the result is that the detour is
worth paying for. Here crossing the critical line carries no additional \emph{singular}
toll --- the cost is finite and generally nonzero, not free --- and the ordered
interior is flat and bounded, while the disordered phase is logarithmically expensive in
$\ln\beta$ as $T\to\infty$. This is the exact opposite of the mean-field antiferromagnet,
where $\beta=0$ was the \emph{cheapest} place in the manifold and the optimum ran to
infinite temperature. Same phenomenon, opposite mechanism; the contrast is worth a sentence
in the paper.

\paragraph{Revised options.}
\begin{enumerate}
  \item \textbf{Keep $H'=0$.} The model delivers the optimal-path result in its sharpest
        form: a theorem rather than a fast-marching computation, in an exactly solved $d=3$
        system. It does not deliver a divergent metric. The scaling analysis of [BR,\S c]
        and Table~I stand on their own and do not require the worked example to exhibit a
        divergence.
  \item \textbf{Restore $H'_\perp$.} Recovers the divergence, at the cost of a
        three-dimensional control space or a slice in which the ordered region is a closed
        curve (Knops Fig.~1) that \emph{can} be circumvented --- which weakens the
        ``crossing is unavoidable'' motivation but sharpens ``around versus through''.
  \item \textbf{Both.} Use the $(\beta,\beta\Hc)$ slice for the optimal-path theorem and a
        $\Hc'_\perp$ slice for the divergence.
\end{enumerate}

Option~3 is probably the strongest paper, and options~1 and 3 both now look better than my
earlier recommendation of option~2. The \textsf{VERIFY} items above still need settling
first, since all of this rests on the finite-metric and flatness results.

\TODO{\S5--\S8 pending the choice above.}

\end{document}
